\begin{center}\begin{tikzpicture}[font=\small,>=stealth]
\draw[black!30,rounded corners=2pt] (0,0) rectangle (14.4,8.2);
\node[anchor=west,font=\small\bfseries] at (.5,7.72) {短轴端点的两条焦半径都等于 $a$，关系就藏在直角三角形里};
\begin{scope}[shift={(7.2,3.8)},scale=.65]
\draw[blue,very thick] (0,0) ellipse (5 and 3);
\draw[black!40] (-5,0)--(5,0);
\draw[orange!85!black,very thick] (-4,0)--(0,3)--(4,0);
\draw[violet,very thick] (0,0)--(0,3);
\draw[black!55] (0,.3)--(.3,.3)--(.3,0);
\fill (-4,0) circle (2pt) node[below] {$F_1$};\fill (4,0) circle (2pt) node[below] {$F_2$};
\fill (0,3) circle (2pt) node[above=6pt] {$B$};
\node[below left] at (0,0) {$O$};
\node[left] at (0,1.5) {$b=3$};
\node[below] at (2,0) {$c=4$};
\node[below left=4pt] at (2,1.5) {$a=5$};
\end{scope}

\node[anchor=west] at (.7,.55) {$BF_1+BF_2=2a$ 且 $BF_1=BF_2$，所以 $BF_2=a$；$a^2=b^2+c^2$};
\end{tikzpicture}\end{center}
